\exercisesection{4}

\begin{enumerate}
\def\labelenumi{\arabic{enumi}.}
\tightlist
\item
  \begin{enumerate}
  \def\labelenumii{(\alph{enumii})}
  \tightlist
  \item
    Let G be a totally ordered group and M a major set in \(G_{+}\) not containing 0. Show that there exists a greatest isolated subgroup H of G not meeting M.
  \end{enumerate}
\end{enumerate}

\begin{enumerate}
\def\labelenumi{(\alph{enumi})}
\setcounter{enumi}{1}
\item
  Let A be a valuation ring and v a valuation associated with A. For an ideal \(p \neq 0\) of A to be prime, it is necessary and sufficient that p correspond to a major set M in the totally ordered value group G of v such that, if H is the greatest isolated subgroup of G not meeting M, M is the complement of \(H_{+}\) in \(G_{+}\) .
\item
  With the notation of (b), for an ideal q of A to be p-primary, it is necessary and sufficient that it satisfy one of the following conditions: either \(H = \{0\}\) (and p is then maximal) or q = p.
\item
  For an ideal \(\mathfrak{a}\) of \(\mathbf{A}\) which is distinct from 0 and \(\mathbf{A}\) to satisfy \(\mathfrak{a}^2 = \mathfrak{a}\) , it is necessary and sufficient that \(\mathfrak{a}\) correspond to a major set \(\mathbf{M}\) such that, if \(\mathbf{H}\) is the greatest isolated subgroup of \(\mathbf{G}\) not meeting \(\mathbf{M}\) , the image \(\mathbf{M}'\) of \(\mathbf{M}\) in the totally ordered group \(\mathbf{G}' = \mathbf{G} / \mathbf{H}\) has no least element; \(\mathfrak{a}\) is then a prime ideal.
\end{enumerate}

\begin{enumerate}
\def\labelenumi{\arabic{enumi}.}
\setcounter{enumi}{1}
\tightlist
\item
  Let G be a totally ordered group.
\end{enumerate}

\begin{enumerate}
\def\labelenumi{(\alph{enumi})}
\item
  For every element a \textgreater{} 0 of G, let \(\mathrm{H}(a)\) be the smallest isolated subgroup containing a, \(\mathrm{H}^{-}(a)\) the greatest isolated subgroup not containing a; show that \(\mathrm{H}(a)/\mathrm{H}^{-}(a)\) is isomorphic to a subgroup of R.
\item
  Conversely, if H is an isolated subgroup of G such that there exists a greatest element (with respect to inclusion) \(H'\) in the set of isolated subgroups of G contained in H and \(\neq H\) , then \(\mathrm{H} = \mathrm{H}(a)\) and \(\mathrm{H}' = \mathrm{H}^{-}(a)\) for all \(a \in H_{+} \cap CH_{+}'\) . Such an isolated subgroup H is called principal and \(H'\) is called its predecessor.
\item
  If \(\mathbf{H}_1, \mathbf{H}_2\) are two distinct isolated subgroups of \(\mathbf{G}\) such that \(\mathbf{H}_1 \subset \mathbf{H}_2\) , show that there exist two isolated subgroups \(\mathbf{H}, \mathbf{H}'\) of \(\mathbf{G}\) such that \(\mathbf{H}_1 \subset \mathbf{H}' \subset \mathbf{H} \subset \mathbf{H}_2\) and \(\mathbf{H} / \mathbf{H}'\) is isomorphic to a non-zero subgroup of \(\mathbf{R}\) .
\item
  Let \(\Sigma\) be the set (totally ordered with respect to inclusion) of isolated subgroups of \(G\) and let \(\Theta \subset \Sigma\) be the set of principal isolated subgroups. Show that \(\Sigma\) is isomorphic to the completion (Set Theory, Chapter III, § 1, Exercise 15) of \(\Theta\) .
\end{enumerate}

\begin{enumerate}
\def\labelenumi{\arabic{enumi}.}
\setcounter{enumi}{2}
\tightlist
\item
  \begin{enumerate}
  \def\labelenumii{(\alph{enumii})}
  \tightlist
  \item
    Let \(\Theta\) be a totally ordered set and for each \(s \in \Theta\) let \(\mathbf{A}_s\) be a subgroup of \(\mathbf{R}\) not reduced to 0. Let
  \end{enumerate}
\end{enumerate}

\[
\Gamma (\Theta , (A _ {s}) _ {s \in \Theta}) = G
\]

be the subgroup of the product \(\biguplus_{s\in\Phi}\mathbf{A}_{s}\) consisting of the functions \(f\) whose support is a well-ordered subset of \(\Theta\) . We define on \(\mathbf{G}\) an ordering compatible with its group structure taking the set \(\mathbf{G}_{+}\) of elements \(\geqslant0\) to be the set consisting of the function 0 and the functions \(f\neq0\) such that \(f(\theta)>0\) for the least element \(\theta\) of the support of \(f\) . Show that \(\mathbf{G}\) is a totally ordered group and that \(\Theta\) is canonically isomorphic to the set of principal isolated subgroups of \(\mathbf{G}\) ordered by the relation \(\supset\) ; moreover, if \(\mathrm{H}(a)\) corresponds to \(s\in\Theta\) under this isomorphism, \(\mathrm{H}(a)/\mathrm{H}^{-}(a)\) is isomorphic to \(\mathbf{A}_{s}\) .

\begin{enumerate}
\def\labelenumi{(\alph{enumi})}
\setcounter{enumi}{1}
\tightlist
\item
  Consider the subgroup \(G'\) of the totally ordered group \(Q \times Q\) (with the lexicographical ordering) generated by the elements \((p_{n}^{-1}, np_{n}^{-1})\) , where \((p_{n})\) is the strictly increasing sequence of prime numbers. Show that the only isolated subgroup \(H'\) of \(G'\) distinct from 0 and \(G'\) is the group \(\{0\} \times Z\) ; but \(G'\) is not isomorphic to a subgroup of the product \(H' \times (G'/H')\) ordered by the lexicographical ordering.
\end{enumerate}

\begin{enumerate}
\def\labelenumi{\arabic{enumi}.}
\setcounter{enumi}{3}
\tightlist
\item
  Let \(A\) be a valuation ring which is not a field.
\end{enumerate}

\begin{enumerate}
\def\labelenumi{(\alph{enumi})}
\item
  Show that for A to be a discrete valuation ring, it is necessary and sufficient that every prime ideal of A be principal.
\item
  Show that, if A is the ring of a valuation of height one, the field of fractions of A is a finitely generated A-algebra.
\end{enumerate}

* 5. Let K be a field and M the set of subrings of K which are not fields. Show that a maximal element of M is a valuation ring of height 1 of its field of fractions L and that K is an algebraic extension of L. The following conditions are equivalent:

(α) M is not empty.

(β) M admits a maximal element.

(γ) There exists a valuation of height 1 on K

(δ) K is not an algebraic extension of a finite prime field (cf.~§ 8, no. 1). *

¶ 6. (a) Let S be a hyperstonian compact space with no isolated point (Integration, Chapter V, § 5, Exercise 14). Let \(\mathcal{C}_{0}(S)\) be the vector space of realvalued functions, finite or otherwise, defined on S, continuous on S and such that the sets \(\overline{f}(-\infty)\) and \(\overline{f}(+\infty)\) are nowhere dense in S (Integration, Chapter II, §1, Exercise 13 (g)). Show that, for every measure \(\mu > 0\) on S, there exist functions f in \(\mathcal{C}_{0}(S)\) whose upper integral is infinite (note that for every nowhere dense closed subset N of S there exists a function f in \(\mathcal{C}_{0}(S)\) which is equal to \(+\infty\) on N; show then that attention may be restricted to the case where the measure \(\mu\) is normal and define f suitably as the upper bound of a sequence of functions in \(\mathcal{C}_{0}(S)\) ).

\begin{enumerate}
\def\labelenumi{(\alph{enumi})}
\setcounter{enumi}{1}
\item
  Let \(\mathbf{G}\) be the ordered subgroup of \(\mathcal{C}_0(S)\) consisting of the functions in \(\mathcal{C}_0(S)\) whose values in \(\mathbf{Z}\) or \(\pm \infty\) . Show that \(\mathbf{G}\) is a complete lattice and is not isomorphic (as an ordered group) to any subgroup of a product group \(\mathbf{R}^{\mathbf{l}}\) .
\item
  Let K be a field and \(A_{0} = K[[X_{s}]]_{s \in S}\) the ring of formal power series with coefficients in K in a family of indeterminates ( \(X_{s}\) ) with the space S as indexing set (Algebra, Chapter IV, § 5, Exercise 1). Let b be the set of formal power series \(P \in A_{0}\) with the following property: there exists a nowhere dense set \(N \subset S\) such that every monomial of P whose coefficient is \(\neq 0\) contains an \(X_{s}\) for which \(s \in N\) . Show that b is an ideal of \(A_{0}\) and that the ring \(A = A_{0}/b\) is an integral domain and completely integrally closed (to prove the latter point, apply Chapter V, § 1, no. 4, Proposition 14, arguing as in Chapter V, § 1, Exercise 10; use also the fact that in S every meagre set is nowhere dense).
\item
  Let \(f\) be an element \(\geqslant 0\) of the group \(G\) and let \(N\) be the nowhere dense set of points where \(f(s) = \pm \infty\) . Let \(p_f\) be the element of \(A\) the image of the formal power series \(\prod_{s \notin N} (1 + X_s)^{f(s)}\) ; the mapping \(f \mapsto p_f\) is an isomorphism of the additive monoid \(G_+\) onto a multiplicative submonoid of \(A\) . Deduce that \(A\) cannot be the intersection of a family of valuation rings of height 1 of its field of fractions (show that this would imply that \(G\) is isomorphic to a subgroup of a product \(R^1\) ). (Cf. Exercises 8 and 9.)
\end{enumerate}

¶ 7. A local integral domain A is said to be of dimension 1 if A is not a field and there is no prime ideal of A distinct from (0) and the maximal ideal m of A. For every ideal a of A, let A: a denote the sub-A-module of the field of fractions of A which is the transporter of a in A (which always contains A). In the following A is assumed to be of dimension 1.

\begin{enumerate}
\def\labelenumi{(\alph{enumi})}
\item
  Show that, if A is strongly Laskerian (Chapter IV, § 2, Exercise 28), then A: m ≠ A. (In the opposite case, note that A: m \(^{r}\) = A for every integer r \textgreater{} 0; on the other hand, every ideal ≠ 0 of A contained in m is m-primary and so in particular is every principal ideal Ax ⊆ m (where x ≠ 0); note finally that m \(^{h}\) ≠ m \(^{k}\) for h ≠ k (Chapter IV, § 2, Exercise 29 (d)).)
\item
  Show that for a strongly Laskerian local integral domain A of dimension 1 the following properties are equivalent:
\end{enumerate}

\((\alpha)\) A is completely integrally closed.

(β) The maximal ideal m is principal.

(γ) Every m-primary ideal is of the form \(m^{k}\) .

(δ) The ideal m is invertible (Chapter II, § 5, no. 6), which is equivalent to m. (A: m) = A.

(ε) A is a discrete valuation ring.

(To show that \((\alpha)\) implies \((\delta)\) , use (a) and observe that if \(\mathfrak{m}.\left(\mathbf{A}:\mathfrak{m}\right)^k = \mathfrak{m}\) for all \(k > 0\) , \(\mathbf{A}\) is not completely integrally closed (cf.~Chapter VII, §1, Exercise 4). To see that \((\delta)\) implies \((\gamma)\) , note that for every ideal \(\mathfrak{a} \neq 0\) contained in \(\mathfrak{m}\) we may write \(\mathfrak{a} = \mathfrak{ma}_1\) , where \(\mathfrak{a}_1 \neq 0\) and use Chapter IV, §2, Exercise 29 (d). Finally, to see that \((\gamma)\) implies \((\beta)\) , observe that \(\mathfrak{m} \neq \mathfrak{m}^2\) and that \((\gamma)\) implies that \(\mathfrak{m} / \mathfrak{m}^2\) is a 1-dimensional vector \((\mathbf{A} / \mathfrak{m})\) -space.)

\begin{enumerate}
\def\labelenumi{(\alph{enumi})}
\setcounter{enumi}{2}
\item
  Show, in the notation of Chapter III, § 3, Exercise 15 (b), that the local ring \(\mathbf{A}_{\mathfrak{m}}\) is a Noetherian integral domain of dimension 1 but is not integrally closed.
\item
  Let K be an algebraically closed field and A the subring of the field \(\mathbf{K}(\mathbf{X},\mathbf{Y})\) of rational functions in two indeterminates over K consisting of the elements \(f\in\mathbf{K}(\mathbf{X},\mathbf{Y})\) such that 0 is substitutable for X in f and \(f(0,\mathbf{Y})\) belongs to K. Show that A is a strongly Laskerian integrally closed local integral domain of dimension 1 but that it is not completely integrally closed. (If B is the polynomial ring \(\mathbf{K}[\mathbf{X},\mathbf{Y}]\) , p the prime ideal BX of B, note that A is the subring of \(B_{p}\) equal to \(K+pB_{p}\) .) (*)
\end{enumerate}

¶ 8. Let A be a local integral domain of dimension 1 (Exercise 7) and K its field of fractions. Show that for A to be the intersection of valuation rings (of K) of height 1, it is necessary and sufficient that A be completely integrally closed. Prove successively that:

\begin{enumerate}
\def\labelenumi{(\alph{enumi})}
\item
  A subring of K containing A and the inverse of an element \(\neq0\) of the maximal ideal m of A is equal to K (use the fact that every ideal of A distinct from 0 and A is m-primary).
\item
  If \(\mathbf{B} \supset \mathbf{A}\) is a subring of \(\mathbf{K}\) distinct from \(\mathbf{K}\) , there exists a valuation ring \(\mathbf{V}\) of height 1 such that \(\mathbf{B} \subset \mathbf{V}\) (use (a)).
\item
  Let \(z \in \mathbf{K} - \mathbf{A}\) ; show that \(zm \subset m\) is impossible arguing as in Exercise 7 (b). Deduce that there exists a valuation ring \(V\) of height 1 such that \(A \subset V\) and \(z \notin V\) . (If \(zm \subset A\) , use (a); otherwise, there is an \(x \in m\) such that
\end{enumerate}

\[
x y = z \in \mathrm{K} - \mathrm{A};
\]

observe that \(z \notin A[z^{-1}]\) and use (b) to show the existence of V.)

¶ 9. Let A be a Laskerian completely integrally closed domain (Chapter IV, §2, Exercise 23). Suppose further, that, for all \(x \neq 0\) in A, the prime ideals \(\mathfrak{p}_i\) weakly associated with A/Ax (Chapter IV, § 1, Exercise 17) are all of height 1, that is to say that for all i, \(p_{i}\) contains no prime ideal other than itself and 0 (cf.~Chapter VII, § 1, no. 6). Show that under these conditions A is an intersection of valuation rings of height 1. (Show first that A is the intersection of all the rings \(A_{p}\) , where p runs the set of prime ideals of height 1, using Exercise 17 (i) of Chapter IV, § 1. Then prove that these local rings \(A_{p}\) are completely integrally closed, using condition ( \(LA_{I}\) ) of Chapter IV, § 2, Exercise 23. Finally apply Exercise 8.) The above hypotheses are in particular satisfied when A is a Laskerian completely integrally closed domain in which every ideal admits a single primary decomposition (Chapter IV, § 2, Exercise 26).

¶ 10. Let A be a ring in which the set of principal ideals is totally ordered by inclusion.

\begin{enumerate}
\def\labelenumi{(\alph{enumi})}
\item
  Show that \(\mathbf{A}\) is a local ring, in which the nilradical \(\mathfrak{N}\) is a prime ideal. Show that, either \(\mathfrak{N}^2 = \mathfrak{N}\) , or \(\mathfrak{N}\) is nilpotent; the ring \(\mathbf{A} / \mathfrak{N}\) is a valuation ring.
\item
  If \(\mathcal{L}\) is the totally ordered set of principal ideals of A, show that the set of ideals of A is totally ordered by inclusion and is isomorphic to the completion of the set \(\mathcal{L}\) (Set Theory, Chapter III, § 1, Exercise 15).
\item
  Suppose that \(\mathfrak{N}^2 = 0\) ; then \(\mathfrak{N}\) may be considered as a module over the valuation ring \(V = A / \mathfrak{N}\) ; let \(K\) be the field of fractions of \(V\) and \(\Gamma\) the order group of a valuation \(v\) on \(K\) corresponding to \(V\) ; show that there exists in \(\Gamma\) two major sets \(M \subset M'\) such that \(\mathfrak{N}\) is a \(V\) -module isomorphic to \(\mathfrak{a}(M') / \mathfrak{a}(M)\) (notation of §3, no. 4).
\item
  Conversely, given a valuation ring V with field of fractions K and order group \(\Gamma\) and two major sets \(M \subset M'\) in \(\Gamma\) , let \(Q = a(M') / a(M)\) and define on the product additive group \(A = V \times Q\) a multiplication by setting
\end{enumerate}

\[
(z, t) (z ^ {\prime}, t ^ {\prime}) = (z z ^ {\prime}, z t ^ {\prime} + z ^ {\prime} t);
\]

show that in A the set of principal ideals is totally ordered by inclusion; the set N of ordered pairs \((0, t)\) where \(t \in Q\) is the nilradical of A, \(N^{2} = 0\) and A/N is isomorphic to V.

\begin{enumerate}
\def\labelenumi{(\alph{enumi})}
\setcounter{enumi}{4}
\item
  For every prime number \(p\) , the ring \(A = \mathbf{Z} / p^2\mathbf{Z}\) is such that the set of principal ideals of \(A\) is totally ordered by inclusion and the nilradical \(\mathfrak{N}\) of \(A\) is such that \(\mathfrak{N}^2 = 0\) ; but show that the ring \(A\) is not isomorphic to the ring constructed (starting with \(V = A / \mathfrak{N}\) and \(Q = \mathfrak{N}\) ) by the method in (d) (note that \(A\) contains no subring isomorphic to \(Z / p\mathbf{Z}\) ).
\item
  Show that for every prime number \(p\) , the algebra \(\mathbf{A}\) of the group \(\mathbf{U}_p\) (Algebra, Chapter VII, §2, Exercise 3) with respect to the prime field \(\mathbf{F}_p\) is a ring in which the principal ideals form a set which is totally ordered by inclusion and the nilradical \(\mathfrak{N}\) satisfies \(\mathfrak{N}^2 = \mathfrak{N}\) , but \(\mathbf{A}\) is not isomorphic to the quotient of a valuation ring by an ideal.
\end{enumerate}

\begin{enumerate}
\def\labelenumi{\arabic{enumi}.}
\setcounter{enumi}{10}
\tightlist
\item
  Let K be a field with a valuation v of height 1.
\end{enumerate}

\begin{enumerate}
\def\labelenumi{(\alph{enumi})}
\tightlist
\item
  Let \(\mathbf{P}(\mathbf{X}) = a_0\mathbf{X}^n +a_1\mathbf{X}^{n - 1} + \dots +a_n\) be a polynomial in \(\mathrm{K[X]}\) of degree \(n > 1\) and such that \(a_{n} \neq 0\) ; show that there exists a strictly increasing sequence \((i_{k})_{0 \leqslant k \leqslant r}\) of integers in the interval \([0, n]\) such that: (1) \(i_{0} = 0\) , \(i_{r} = n\) ; (2) \(v(a_{i_{k}})\) is finite for \(0 \leqslant k \leqslant r\) ; (3) for every index \(j\) such that \(0 \leqslant j \leqslant n\) distinct from the \(i_{k}\) , such that \(v(a_{j})\) is finite, the point
\end{enumerate}

\[
(j, v (a _ {j})) \in \mathbf {R} ^ {2}
\]

lies above the line passing through the points \((i_{k}, v(a_{i_{k}}))\) and \((i_{k+1}, v(a_{i_{k+1}}))\) and strictly above that line if \(j < i_{k}\) or \(j > i_{k+1}\) . The union of the segments joining the points \((i_{k}, v(a_{i_{k}}))\) and \((i_{k+1}, v(a_{i_{k+1}}))\) is called the Newton polygon of P, the above segments are called the sides and the points \((i_{k}, v(a_{i_{k}}))\) the vertices of the polygon.

\begin{enumerate}
\def\labelenumi{(\alph{enumi})}
\setcounter{enumi}{1}
\item
  Suppose that all the zeros of P belong to K. Show that for the valuations of all the zeros of P to be the same, it is necessary and sufficient that r = 1 (in other words, that the Newton polygon reduce to a single side). (To show that the condition is sufficient, consider the Newton polygon of a product \(P_{1}P_{2}\) where all the zeros of \(P_{1}\) are invertible in the ring of v, whilst all those of \(P_{2}\) belong to the ideal of v.)
\item
  Suppose that all the zeros of \(\mathbf{P}\) belong to \(\mathbf{K}\) ; form the Newton polygon of \(\mathbf{P}\) and write.
\end{enumerate}

\[
\rho_ {k} = i _ {k + 1} - i _ {k}, \quad \sigma_ {k} = (v (a _ {i _ {k + 1}}) - v (a _ {i _ {k}})) / \rho_ {k}.
\]

Show that, for \(0 \leqslant k \leqslant r - 1\) , P admits exactly \(\rho_k\) zeros (counted with their orders of multiplicity) whose valuations are all equal to \(\sigma_k\) (use (b) and argue by induction on \(r\) ).

\begin{enumerate}
\def\labelenumi{(\alph{enumi})}
\setcounter{enumi}{3}
\tightlist
\item
  Generalize to the case of any valuation v (embed the order group \(\Gamma\) of v in the vector Q-space \(\Gamma_{(Q)}\) , which has a natural totally ordered group structure).
\end{enumerate}

